\begin{titlepage}
\thispagestyle{empty}
\centering
\vspace*{1.2cm}
{\large UNIVERSIDAD NACIONAL DE COLOMBIA\par}
{\large Sede La Paz\par}
\vspace{0.3cm}
{\normalsize Ingeniería Mecatrónica --- Proyecto Académico Especial (PAE)\par}
{\normalsize Semillero de Investigación PRIMBIO --- Grupo de Investigación IAM-NANO\par}
\vspace{2.6cm}
{\color{azul}\rule{0.85\textwidth}{1.2pt}\par}
\vspace{0.5cm}
{\Huge\bfseries Modelo Matemático del Simulador\\[0.25cm] QuadSIM-PRIMBIO\par}
\vspace{0.4cm}
{\LARGE 3.ª edición\par}
\vspace{0.4cm}
{\color{azul}\rule{0.85\textwidth}{1.2pt}\par}
\vspace{1.2cm}
{\large Dinámica no lineal de seis grados de libertad, propulsión, sensores,\\
estimación, control en cascada y controladores de trayectoria,\\
con verificación numérica de cada bloque\par}
\vfill
{\Large Jesús Daniel Jácome\par}
\vspace{0.25cm}
{\large Director: Francisco Javier Burgos Flórez\par}
\vspace{0.6cm}
{\large Universidad Nacional de Colombia, Sede La Paz\par}
{\large 2026\par}
\vspace{1.2cm}
\begin{minipage}{0.86\textwidth}\small\color{tinta2}
Documento de referencia del paquete \cod{quadsim/} de QuadSIM-PRIMBIO 3.0
(repositorio \cod{PAE-Quadcopter-SIM}). Todas las cifras marcadas como evidencia numérica
salen de \cod{experiments/model\_verification.py}; el archivo
\cod{docs/modelo/valores.tex} se regenera a partir de sus resultados en cada
compilación, de modo que el texto no puede quedar desfasado del código.
\end{minipage}
\end{titlepage}
